% Generated by publication/build_sections.py; do not edit by hand.
\paragraph{Informal verdict.}
\textbf{C --- Substantial progress}, with high confidence in the stated partial results.

\paragraph{Lean coverage (partial).}
Verified Lean covers the 667-line intrinsic-distance/chord core, generic cone lower bounds, Clifford coordinates, and the exact all-pairs singular 2-by-2 formula across seven imported modules (3,143 lines). It does not prove generic cone upper equalities, equality with the PDF finite piecewise-C1 distance, complete positive chord/boundary theory, arbitrary-n generic distances in GL\_n\textasciicircum{}+ or Sigma\_n, or the canonical attainment classification.

\subsection{Audited informal answer}
\subsubsection{Faithful restatement}

Equip \(M_n(\mathbb R)\cong \mathbb R^{n^2}\) with its standard Euclidean structure, hence with the Frobenius inner product and norm
\[
\langle X,Y\rangle_F=\operatorname{tr}(X^{\mathsf T}Y),
\qquad
|X|*F=\sqrt{\operatorname{tr}(X^{\mathsf T}X)}.
\]
For a piecewise \(C^1\)-arc-connected subset \(E\subset M_n(\mathbb R)\), its intrinsic length distance is
\[
d_E(A,B)=\inf*\gamma \int_0^1|\gamma'(t)|_F,dt,
\]
where \(\gamma\) ranges over piecewise \(C^1\) curves in \(E\) joining \(A\) to \(B\).

The problem asks for this distance first on
\[
\Omega_n:=GL_n^+(\mathbb R)={X:\det X>0},
\]
and then on
\[
\Sigma_n:={X:\det X=0},
\]
the set called ``matrices non régulières'' in the PDF. The duplicated phrase ``\(x\) et \(y\)'' in the first sentence is only a typographical error; there is no mathematical ambiguity. ([Luc Pommeret][1])

\medskip\hrule\medskip

\subsubsection{Status}

\textbf{Partial result.} I give:

\begin{enumerate}
\item an exact metric-cone reduction valid for every \(n\);
\item a complete characterization of all pairs for which the answer is simply the Frobenius distance;
\item quantitative strict lower bounds when it is not;
\item explicit constructive upper bounds;
\item a closed formula for all pairs of singular \(2\times2\) matrices;
\item an exact scalar variational determination for \(GL_2^+(\mathbb R)\).
\end{enumerate}

What remains unresolved here is an explicit evaluation, for arbitrary \(n\ge 3\), of certain intrinsic angular distances on the unit determinant links. This is the exact obstruction, not merely a technical gap in one estimate.

\medskip\hrule\medskip

\subsection{1. Exact reduction to angular distances on the unit sphere}

Let
\[
S_F={X\in M_n(\mathbb R):|X|*F=1},
\]
and define the links
\[
Y_n^+:=\Omega_n\cap S_F,
\qquad
Z_n:=\Sigma_n\cap S_F.
\]
Write \(\delta*+\) and \(\delta_\Sigma\) for their intrinsic spherical length distances. Thus curve lengths are still computed with the ambient Frobenius metric, restricted to \(S_F\).

For a nonzero matrix \(A\), put
\[
r_A=|A|_F,\qquad \widehat A=\frac{A}{|A|_F}.
\]

\subsubsection{Theorem 1 --- Metric-cone formula}

For \(A,B\in\Omega_n\),
\[
\boxed{
d_{\Omega_n}(A,B)^2
=
%
r_A^2+r_B^2
-2r_Ar_B
\cos\!\left(\min{\pi,\delta_+(\widehat A,\widehat B)}\right).
}
\tag{1}
\]

For nonzero \(A,B\in\Sigma_n\),
\[
\boxed{
d_{\Sigma_n}(A,B)^2
=
%
r_A^2+r_B^2
-2r_Ar_B
\cos\!\left(\min{\pi,\delta_\Sigma(\widehat A,\widehat B)}\right).
}
\tag{2}
\]
If \(A=0\) or \(B=0\), then
\[
\boxed{d_{\Sigma_n}(A,B)=|A-B|_F.}
\tag{3}
\]

Consequently, in both sets,
\[
\boxed{
|A-B|_F\le d(A,B)\le |A|_F+|B|_F.
}
\tag{4}
\]

The upper bound in \(\Sigma_n\) is attained by the broken radial path \(A\to0\to B\). In \(GL_n^+\), it is in general only an infimum, obtained by passing arbitrarily close to \(0\).

\medskip\hrule\medskip

\subsubsection{Proof of the metric-cone formula}

We prove a general lemma.

\paragraph{Lemma 2 --- Distance in a Euclidean cone}

Let \(Y\subset S^{N-1}\) be piecewise \(C^1\)-arc-connected, and let
\[
C^\times(Y)={ru:r>0,\ u\in Y}.
\]
Denote by \(\delta_Y\) the intrinsic spherical distance on \(Y\). For \(x=ru\) and \(y=sv\),
\[
d_{C^\times(Y)}(x,y)^2
=
%
r^2+s^2-2rs\cos\!\left(\min{\pi,\delta_Y(u,v)}\right).
\tag{5}
\]

The same formula holds for the closed cone
\[
C(Y)={ru:r\ge0,\ u\in Y},
\]
with the evident convention when an endpoint is \(0\).

\paragraph{Proof}

Let \(\gamma(t)=\rho(t)u(t)\) be a piecewise \(C^1\) curve in the punctured cone, where \(\rho(t)>0\) and \(u(t)\in Y\). Since \(|u(t)|=1\),
\[
\langle u,u'\rangle=0,
\]
and hence
\[
|\gamma'(t)|^2
=
%
\rho'(t)^2+\rho(t)^2|u'(t)|^2.
\tag{6}
\]

Define the accumulated angular length
\[
q(t)=\int_0^t|u'(\tau)|,d\tau
\]
and the developed planar curve
\[
\widetilde\gamma(t)
=
%
\rho(t)\bigl(\cos q(t),\sin q(t)\bigr)\in\mathbb R^2.
\]
By \(6\),
\[
|\widetilde\gamma'(t)|=|\gamma'(t)|
\]
almost everywhere. Thus the two curves have the same length.

Let \(L=q(1)\). Necessarily
\[
L\ge \delta_Y(u,v).
\]

If \(L\le\pi\), the planar chord inequality gives
\[
\operatorname{length}(\gamma)
\ge
\sqrt{r^2+s^2-2rs\cos L}.
\]
The right-hand side is increasing in \(L\in[0,\pi]\), so
\[
\operatorname{length}(\gamma)
\ge
\sqrt{r^2+s^2-2rs\cos\delta_Y(u,v)}
\]
whenever \(\delta_Y(u,v)\le L\le\pi\).

If \(L\ge\pi\), let \(t_0\) be the first time at which \(q(t_0)=\pi\), and put \(a=\rho(t_0)\). In the developed plane, \(\widetilde\gamma(t_0)=(-a,0)\). Therefore
\[
\operatorname{length}(\gamma|_{[0,t_0]})\ge r+a,
\]
while, by the reverse triangle inequality,
\[
\operatorname{length}(\gamma|_{[t_0,1]})\ge |s-a|.
\]
Hence
\[
\operatorname{length}(\gamma)
\ge r+a+|s-a|
\ge r+s.
\]
This proves the lower bound in \(5\).

For the reverse inequality, first suppose
\[
\delta_Y(u,v)<\pi.
\]
Choose a unit-speed curve
\[
\sigma:[0,L]\longrightarrow Y,
\qquad
\sigma(0)=u,\quad \sigma(L)=v,
\]
with \(L<\pi\) and \(L\) arbitrarily close to \(\delta_Y(u,v)\).

The straight segment in \(\mathbb R^2\) joining
\[
r(1,0)
\quad\text{to}\quad
s(\cos L,\sin L)
\]
can be written in polar form
\[
\theta\longmapsto
\rho(\theta)(\cos\theta,\sin\theta),
\qquad 0\le\theta\le L,
\]
where
\[
\rho(\theta)
=
%
\frac{rs\sin L}
{s\sin(L-\theta)+r\sin\theta}>0.
\]
Lift this planar segment to the cone by
\[
\Gamma(\theta)=\rho(\theta)\sigma(\theta).
\]
Its length equals the length of the planar chord:
\[
\sqrt{r^2+s^2-2rs\cos L}.
\]
Letting \(L\downarrow\delta_Y(u,v)\) proves the upper bound.

If \(\delta_Y(u,v)\ge\pi\), fix any finite-length curve in \(Y\) from \(u\) to \(v\), of length \(L\). For \(\varepsilon>0\), concatenate:

\[
ru\longrightarrow \varepsilon u,
\qquad
\varepsilon u\longrightarrow \varepsilon v,
\qquad
\varepsilon v\longrightarrow sv.
\]
Its length is
\[
(r-\varepsilon)+\varepsilon L+(s-\varepsilon)\longrightarrow r+s.
\]
This proves \(5\). If the apex belongs to the cone, the path through \(0\) gives length exactly \(r+s\). \(\square\)

Applying this lemma to \(Y_n^+\) and \(Z_n\) proves Theorem 1.

For completeness, \(Y_n^+\) is piecewise \(C^1\)-arc-connected by polar decomposition. Indeed,
\[
A=QP,\qquad Q\in SO(n),\quad P>0,
\]
and one connects \(P\) to \(I\) through positive-definite matrices, then \(Q\) to \(I\) inside \(SO(n)\).

For \(n\ge2\), \(Z_n\) is also piecewise \(C^1\)-arc-connected. A nonzero singular matrix can, using a singular-value decomposition, be deformed on the unit sphere to a unit rank-one matrix. The unit rank-one locus is the continuous image of
\[
S^{n-1}\times S^{n-1},\qquad (u,v)\longmapsto uv^{\mathsf T},
\]
and is therefore connected.

\medskip\hrule\medskip

\subsection{2. Exactly when the distance is the Frobenius chord}

The answer is frequently, but not always, \(|A-B|_F\). There is a complete algebraic criterion.

\subsubsection{Theorem 3 --- Chord criterion in \(GL_n^+\)}

For \(A,B\in GL_n^+(\mathbb R)\), the following are equivalent:

\[
\boxed{d_{\Omega_n}(A,B)=|A-B|_F;}
\tag{7}
\]

\[
\boxed{\det((1-t)A+tB)\ge0\quad\text{for every }t\in[0,1];}
\tag{8}
\]

\[
\boxed{\text{every negative real eigenvalue of }A^{-1}B
\text{ has even algebraic multiplicity}.}
\tag{9}
\]

Notice that condition \(8\) allows the segment to touch the singular locus. Thus the infimum can equal the Euclidean chord even though the straight segment itself is not an admissible curve in \(GL_n^+\).

\medskip\hrule\medskip

\subsubsection{Theorem 4 --- Chord criterion in the singular locus}

For \(A,B\in\Sigma_n\), the following are equivalent:

\[
\boxed{d_{\Sigma_n}(A,B)=|A-B|_F;}
\tag{10}
\]

\[
\boxed{\det((1-t)A+tB)=0\quad\text{for every }t\in[0,1];}
\tag{11}
\]

\[
\boxed{\det(A+\lambda B)\equiv0\quad\text{as a polynomial in }\lambda.}
\tag{12}
\]

Equivalently, the pencil \(A+\lambda B\) is singular over \(\mathbb R(\lambda)\). This is equivalent to the existence of a nonzero polynomial vector
\[
v(\lambda)\in\mathbb R[\lambda]^n
\]
such that
\[
(A+\lambda B)v(\lambda)=0.
\tag{13}
\]

A common right or left null vector is sufficient for \(11\), but is not necessary.

The proofs use two elementary geometric lemmas.

\medskip\hrule\medskip

\subsubsection{2.1 A quantitative obstacle lemma}

\paragraph{Lemma 5}

Let \(x,y\in\mathbb R^N\), let \(D=|x-y|\), and set
\[
c_t=(1-t)x+ty,\qquad 0<t<1.
\]
Suppose a continuous curve \(\gamma\) from \(x\) to \(y\) avoids the open ball
\[
B(c_t,\rho).
\]
Then
\[
\boxed{
\operatorname{length}(\gamma)
\ge
\sqrt{t^2D^2+\rho^2}
+
\sqrt{(1-t)^2D^2+\rho^2}.
}
\tag{14}
\]

\paragraph{Proof}

Let
\[
e=\frac{y-x}{D}.
\]
The continuous function
\[
s\longmapsto \langle\gamma(s)-x,e\rangle
\]
takes the values \(0\) and \(D\). Thus at some point \(p\) on the curve,
\[
\langle p-x,e\rangle=tD.
\]
Write
\[
p=c_t+v,\qquad v\perp e.
\]
Since the curve avoids \(B(c_t,\rho)\),
\[
|v|\ge\rho.
\]
The two subarcs have lengths at least their chord lengths, so
\[
\operatorname{length}(\gamma)
\ge
|p-x|+|y-p|,
\]
which gives \(14\). \(\square\)

\medskip\hrule\medskip

\subsubsection{2.2 Distance to the singular locus}

For an invertible matrix \(C\), let \(\sigma_{\min}(C)\) denote its least singular value.

\paragraph{Lemma 6}

\[
\boxed{
\operatorname{dist}*F(C,\Sigma_n)=\sigma_{\min}(C).
}
\tag{15}
\]

If \(\det C<0\), then also
\[
\boxed{
\operatorname{dist}*F(C,{\det\ge0})=\sigma_{\min}(C).
}
\tag{16}
\]

\paragraph{Proof}

Let \(S\) be singular and choose a unit vector \(v\in\ker S\). Then
\[
|C-S|_F
\ge
|(C-S)v|
=
%
|Cv|
\ge
\sigma_{\min}(C).
\]
Conversely, let \(Cv=\sigma_{\min}(C)u\) be a least-singular-value pair. Then
\[
S=C-\sigma_{\min}(C)uv^{\mathsf T}
\]
satisfies \(Sv=0\) and
\[
|C-S|*F=\sigma_{\min}(C).
\]
This proves \(15\).

If \(\det C<0\) and \(P\) has nonnegative determinant, the segment from \(C\) to \(P\) meets \(\Sigma_n\). Hence
\[
|C-P|_F\ge\operatorname{dist}_F(C,\Sigma_n).
\]
Since \(\Sigma_n\subset{\det\ge0}\), equality of the two distances follows. \(\square\)

Combining Lemmas 5 and 6 yields the following useful strict estimates.

\subsubsection{Corollary 7 --- Quantitative strictness}

Let
\[
C_t=(1-t)A+tB,\qquad D=|A-B|_F.
\]

If \(A,B\in GL_n^+\) and \(\det C_t<0\), then
\[
\boxed{
d_{\Omega_n}(A,B)
\ge
\sqrt{t^2D^2+\sigma_{\min}(C_t)^2}
+
\sqrt{(1-t)^2D^2+\sigma_{\min}(C_t)^2}
%
D.
}
\tag{17}
\]

If \(A,B\in\Sigma_n\) and \(C_t\) is invertible, then
\[
\boxed{
d_{\Sigma_n}(A,B)
\ge
\sqrt{t^2D^2+\sigma_{\min}(C_t)^2}
+
\sqrt{(1-t)^2D^2+\sigma_{\min}(C_t)^2}
%
D.
}
\tag{18}
\]

The lower bounds may of course be maximized over all admissible \(t\).

\medskip\hrule\medskip

\subsubsection{2.3 A local bypass lemma}

To prove the sufficiency of \(8\), one must show that isolated contacts of the segment with \(\Sigma_n\) can be bypassed at arbitrarily small cost.

\paragraph{Lemma 8 --- Local positive-determinant bypass}

For each \(C\in\Sigma_n\), there exist constants \(h_0,L,R>0\) such that, whenever \(0<h<h_0\) and
\[
X,Y\in GL_n^+\cap B(C,h),
\]
there is a piecewise \(C^1\) path in
\[
GL_n^+\cap B(C,Rh)
\]
joining \(X\) to \(Y\), of length at most \(Lh\).

\paragraph{Proof}

Let \(k=\operatorname{rank}C\), \(m=n-k\). By left and right orthogonal transformations that preserve Frobenius distances and determinant sign, we may suppose
\[
C=
\begin{pmatrix}
D&0\\
0&0
\end{pmatrix},
\]
where \(D\) is a positive diagonal \(k\times k\) matrix.

The determinant-sign-preserving reduction deserves one detail. Starting with a singular-value decomposition, if the product of the determinants of the two orthogonal factors is (-1), one flips one basis vector in a zero singular direction. This does not alter \(C\), but changes that product to (+1).

Write matrices near \(C\) in block form
\[
M=
\begin{pmatrix}
P&F\\
G&H
\end{pmatrix}.
\]
For \(M\) sufficiently close to \(C\), \(P\) is invertible and \(\det P>0\). Introduce the Schur complement
\[
W=H-GP^{-1}F.
\]
The map
\[
M\longmapsto(P,F,G,W)
\]
is a smooth local diffeomorphism with inverse
\[
(P,F,G,W)\longmapsto
\begin{pmatrix}
P&F\\
G&W+GP^{-1}F
\end{pmatrix}.
\]
On a sufficiently small fixed neighborhood, this map and its inverse are bi-Lipschitz. Moreover,
\[
\det M=\det P,\det W,
\]
so in these coordinates
\[
\det M>0\iff \det W>0.
\]

It remains to connect two small matrices \(W_0,W_1\in GL_m^+\) by a path of size and length \(O(h)\). Write their polar decompositions
\[
W_i=Q_iH_i,\qquad Q_i\in SO(m),\quad H_i>0.
\]
Fix \(a=h\). Connect \(W_i\) to \(aQ_i\) through
\[
Q_i\bigl((1-s)H_i+saI\bigr).
\]
This stays in \(GL_m^+\) and has length
\[
|H_i-aI|_F=O(h).
\]

Next connect \(aQ_0\) to \(aQ_1\) in \(aSO(m)\). If
\[
Q_0^{\mathsf T}Q_1=e^K,\qquad K^{\mathsf T}=-K,
\]
then
\[
s\longmapsto aQ_0e^{sK}
\]
has length \(a|K|_F=O(h)\). Such a skew-symmetric logarithm always exists for an element of \(SO(m)\); one may take the usual block logarithm on its planar rotation blocks, with all angles in \([-\pi,\pi]\).

Interpolate the free coordinates (P,F,G) linearly while following this path in \(W\). The bi-Lipschitz inverse coordinate map gives the asserted path and estimates. The case \(k=0\) is the same argument with the (P,F,G) blocks absent. \(\square\)

\medskip\hrule\medskip

\subsubsection{Proof of Theorem 3}

Set
\[
p(t)=\det((1-t)A+tB).
\]

If \(p(t)<0\) for some \(t\), Corollary 7 gives
\[
d_{\Omega_n}(A,B)>|A-B|_F.
\]
Thus \(7\) implies \(8\).

Conversely, suppose \(p(t)\ge0\) on ([0,1]). Since \(p(0),p(1)>0\), \(p\) is not the zero polynomial and has only finitely many zeros
\[
0<t_1<\cdots<t_N<1.
\]
Outside arbitrarily small disjoint intervals around the \(t_j\), the straight segment lies in \(GL_n^+\).

Inside an interval around \(t_j\), its two endpoints lie in \(GL_n^+\) and tend to the singular matrix
\[
C_j=(1-t_j)A+t_jB.
\]
Lemma 8 replaces that short portion by a path in \(GL_n^+\) whose length tends to \(0\) with the interval size. Replacing all such portions gives curves in \(GL_n^+\) whose lengths tend to
\[
|A-B|*F.
\]
Since every curve has length at least its Euclidean chord,
\[
d_{\Omega_n}(A,B)=|A-B|_F.
\]
This proves the equivalence of \(7\) and \(8\).

To prove the spectral formulation, put
\[
C=A^{-1}B.
\]
Then
\[
p(t)=\det A\cdot \det((1-t)I+tC).
\]
Over \(\mathbb C\),
\[
\det((1-t)I+tC)
=
%
\prod_{\lambda\in\sigma(C)}
(1-t+t\lambda),
\]
with eigenvalues repeated according to algebraic multiplicity.

A positive real eigenvalue contributes a positive factor on ([0,1]). A nonreal conjugate pair \(\lambda,\bar\lambda\) contributes
\[
|1-t+t\lambda|^2>0.
\]
A negative real eigenvalue \(\lambda\) contributes a zero at
\[
t=\frac1{1-\lambda}\in(0,1),
\]
whose multiplicity is precisely the algebraic multiplicity of \(\lambda\). Hence the determinant is nonnegative throughout ([0,1]) exactly when each such multiplicity is even. \(\square\)

\medskip\hrule\medskip

\subsubsection{Proof of Theorem 4}

If
\[
\det((1-t)A+tB)\equiv0,
\]
the straight segment lies in \(\Sigma_n\), and therefore
\[
d_{\Sigma_n}(A,B)=|A-B|_F.
\]

If the determinant polynomial is not identically zero, there is some \(t\in(0,1)\) for which \(C_t\) is invertible. Corollary 7 then gives strict inequality. This proves \(10\)--\(11\).

The equivalence with \(12\) follows by homogenizing the determinant polynomial. Finally, a square matrix over the field \(\mathbb R(\lambda)\) has zero determinant exactly when it has a nonzero kernel vector over that field. Clearing denominators gives \(13\). \(\square\)

\medskip\hrule\medskip

\subsection{3. Constructive upper bounds in arbitrary dimension}

The cone formula is exact but leaves an angular metric. The following bounds give concrete finite-dimensional constructions.

\subsubsection{3.1 A polar-decomposition route in \(GL_n^+\)}

Let
\[
A=Q_AP_A,\qquad B=Q_BP_B,
\]
be the polar decompositions, with
\[
Q_A,Q_B\in SO(n),\qquad P_A,P_B>0.
\]
Let the eigenvalues of \(Q_A^{\mathsf T}Q_B\) be
\[
e^{\pm i\theta_1},\ldots,e^{\pm i\theta_m},
\qquad 0\le\theta_j\le\pi,
\]
together with possible eigenvalues \(1\), and put
\[
\vartheta(Q_A,Q_B)
=
%
\sqrt{2\sum_{j=1}^m\theta_j^2}.
\]
This is the length of the standard shortest rotation path in \(SO(n)\) for the induced Frobenius metric.

For every \(\varepsilon>0\), the concatenation
\[
A\longrightarrow \varepsilon Q_A
\longrightarrow \varepsilon Q_B
\longrightarrow B
\]
stays in \(GL_n^+\), and gives
\[
\boxed{
d_{\Omega_n}(A,B)
\le
\inf_{\varepsilon>0}
\left(
|P_A-\varepsilon I|_F
+
\varepsilon\vartheta(Q_A,Q_B)
+
|P_B-\varepsilon I|_F
\right).
}
\tag{19}
\]

Taking \(\varepsilon\downarrow0\) recovers the bound
\[
d_{\Omega_n}(A,B)\le|A|_F+|B|_F.
\]

\paragraph{Exact special case: distance to \(SO(n)\)}

The polar factor is the unique closest element of \(SO(n)\), even for the intrinsic distance:
\[
\boxed{
\operatorname{dist}_{d_{\Omega_n}}(A,SO(n))
=
%
|P_A-I|_F,
}
\tag{20}
\]
and the unique minimizer is \(Q_A\).

Indeed, the segment
\[
Q_A\bigl((1-t)P_A+tI\bigr)
\]
lies in \(GL_n^+\) and has length \(|P_A-I|_F\). Conversely, intrinsic distance dominates Euclidean distance, and \(Q_A\) is the unique Frobenius-nearest rotation.

\medskip\hrule\medskip

\subsubsection{3.2 A common-kernel construction in the singular locus}

For a unit vector \(x\), put
\[
P_x=I-xx^{\mathsf T}.
\]
Then \(AP_x\) and \(BP_x\) both annihilate \(x\). The three paths
\[
A\longrightarrow AP_x,
\qquad
AP_x\longrightarrow BP_x,
\qquad
BP_x\longrightarrow B
\]
all lie in \(\Sigma_n\). Their total length is
\[
|Ax|
+
|(A-B)P_x|*F
+
|Bx|.
\]
Therefore
\[
\boxed{
d_{\Sigma_n}(A,B)
\le
\min_{|x|=1}
\left(
|Ax|
+
|(A-B)(I-xx^{\mathsf T})|_F
+
|Bx|
\right).
}
\tag{21}
\]

By applying this to transposes, one also gets
\[
\boxed{
d_{\Sigma_n}(A,B)
\le
\min_{|y|=1}
\left(
|A^{\mathsf T}y|
+
|(I-yy^{\mathsf T})(A-B)|_F
+
|B^{\mathsf T}y|
\right).
}
\tag{22}
\]

If \(A\) and \(B\) have a common right or left null vector, these bounds give the Euclidean chord exactly.

\medskip\hrule\medskip

\subsection{4. Complete solution for singular \(2\times2\) matrices}

For
\[
A=
\begin{pmatrix}
a&b\ c&d
\end{pmatrix},
\]
define
\[
\Phi(A)=(z,w)\in\mathbb C^2
\]
by
\[
z=\frac{a+d+i(c-b)}{\sqrt2},
\qquad
w=\frac{a-d+i(c+b)}{\sqrt2}.
\tag{23}
\]

A direct computation gives
\[
|z|^2+|w|^2=|A|_F^2,
\tag{24}
\]
and
\[
\det A=\frac{|z|^2-|w|^2}{2}.
\tag{25}
\]
Thus \(\Phi\) is a real-linear isometry, and
\[
\Sigma_2
\quad\longleftrightarrow\quad
{(z,w):|z|=|w|}.
\]

The unit link is therefore the Clifford torus
\[
T=
\left\{
\frac1{\sqrt2}(e^{i\alpha},e^{i\beta})
:
\alpha,\beta\in\mathbb R
\right\}.
\]
Its induced metric is
\[
ds_T^2=\frac12(d\alpha^2+d\beta^2).
\tag{26}
\]

For nonzero complex numbers \(\zeta,\xi\), define their principal angular separation by
\[
\operatorname{ang}(\zeta,\xi)
=
%
\min_{k\in\mathbb Z}
\bigl|
\arg\xi-\arg\zeta+2\pi k
\bigr|
\in[0,\pi].
\]

\subsubsection{Theorem 9 --- Exact distance on \(\Sigma_2\)}

Let \(A,B\in\Sigma_2\setminus{0}\), with
\[
\Phi(A)=(z_A,w_A),
\qquad
\Phi(B)=(z_B,w_B).
\]
Define
\[
\Theta(A,B)
=
%
\sqrt{
\frac{
\operatorname{ang}(z_A,z_B)^2
+
\operatorname{ang}(w_A,w_B)^2
}{2}
}.
\tag{27}
\]
Then \(0\le\Theta(A,B)\le\pi\), and
\[
\boxed{
d_{\Sigma_2}(A,B)
=
%
\sqrt{
|A|_F^2+|B|_F^2
---------------
%
2|A|_F|B|_F\cos\Theta(A,B)
}.
}
\tag{28}
\]

If one endpoint is \(0\), the distance is the Frobenius norm of the other.

\paragraph{Proof}

The universal cover of \(T\) is \(\mathbb R^2\) with metric
\[
\frac12(d\alpha^2+d\beta^2).
\]
Hence the intrinsic torus distance between the normalized directions of \(A\) and \(B\) is exactly \(\Theta(A,B)\). Since \(\Theta\le\pi\), formula \(28\) follows directly from the metric-cone formula. \(\square\)

\paragraph{Example}

Let
\[
A=E_{11},
\qquad
B=E_{22}.
\]
Then
\[
\Phi(A)=\frac1{\sqrt2}(1,1),
\qquad
\Phi(B)=\frac1{\sqrt2}(1,-1).
\]
Thus
\[
\Theta(A,B)=\frac{\pi}{\sqrt2},
\]
and
\[
\boxed{
d_{\Sigma_2}(E_{11},E_{22})
=
%
\sqrt{2-2\cos(\pi/\sqrt2)}
\approx1.79204.
}
\]
This is strictly greater than
\[
|E_{11}-E_{22}|_F=\sqrt2.
\]

At the midpoint, \(\frac12(E_{11}+E_{22})=\frac12I\) is invertible. The general obstacle bound already gives
\[
d_{\Sigma_2}(E_{11},E_{22})\ge\sqrt3,
\]
while \(28\) gives the exact value.

\medskip\hrule\medskip

\subsection{5. Exact variational determination for \(GL_2^+(\mathbb R)\)}

Under the same isometry \(\Phi\),
\[
GL_2^+(\mathbb R)
\quad\longleftrightarrow\quad
{(z,w)\in\mathbb C^2:|z|>|w|}.
\]
Its unit link is the open solid Clifford torus
\[
H^\circ
=
%
{(z,w)\in S^3:|z|>|w|}.
\]

Use Hopf coordinates
\[
z=\cos\eta,e^{i\alpha},
\qquad
w=\sin\eta,e^{i\beta},
\qquad
0\le\eta<\frac{\pi}{4}.
\tag{29}
\]
The round metric on \(S^3\) is
\[
ds^2
=
%
d\eta^2+\cos^2\eta,d\alpha^2+\sin^2\eta,d\beta^2.
\tag{30}
\]

Let \(p,q\in H^\circ\) have coordinates
\[
(\eta_0,\alpha_0,\beta_0),
\qquad
(\eta_1,\alpha_1,\beta_1).
\]
Let
\[
a=\min_{k\in\mathbb Z}
|\alpha_1-\alpha_0+2\pi k|\in[0,\pi].
\]
If both \(w\)-coordinates are nonzero, let
\[
b=\min_{\ell\in\mathbb Z}
|\beta_1-\beta_0+2\pi\ell|\in[0,\pi].
\]
If one \(w\)-coordinate is zero, its \(\beta\)-coordinate is immaterial, and we set \(b=0\).

For
\[
\eta\in H^1([0,1]),
\qquad
\eta(0)=\eta_0,\quad\eta(1)=\eta_1,
\qquad
0\le\eta(t)\le\frac{\pi}{4},
\]
define
\[
I_c(\eta)=\int_0^1\sec^2\eta(t),dt,
\qquad
I_s(\eta)=\int_0^1\csc^2\eta(t),dt,
\tag{31}
\]
allowing \(I_s(\eta)=+\infty\), and using the convention
\[
\frac{b^2}{+\infty}=0.
\]

\subsubsection{Theorem 10 --- Scalar variational formula for the link}

The intrinsic spherical distance in \(H^\circ\) satisfies
\[
\boxed{
\delta_H(p,q)^2
=
%
\inf_{\eta}
\left[
\int_0^1\eta'(t)^2,dt
+
\frac{a^2}{I_c(\eta)}
+
\frac{b^2}{I_s(\eta)}
\right].
}
\tag{32}
\]

Consequently, for \(A,B\in GL_2^+\), with
\[
r=|A|_F,\quad s=|B|*F,\qquad
p=\frac{\Phi(A)}r,\quad q=\frac{\Phi(B)}s,
\]
one has the exact formula
\[
\boxed{
d_{\Omega_2}(A,B)
=
%
\sqrt{
r^2+s^2
-------
%
2rs\cos\!\left(\min{\pi,\delta_H(p,q)}\right)
}.
}
\tag{33}
\]

Thus the \(2\times2\) problem reduces exactly to a one-variable obstacle variational problem.

\paragraph{Proof of \(32\)}

For a unit-time curve in \(H^\circ\), the energy is
\[
E=
\int_0^1
\left(
\eta'^2+\cos^2\eta,\alpha'^2+\sin^2\eta,\beta'^2
\right)dt.
\tag{34}
\]
The infimum of \(E\) is the square of the length distance: every curve can be reparameterized at constant speed, and Cauchy--Schwarz gives \(L^2\le E\).

Fix \(\eta\). By Cauchy--Schwarz,
\[
\left(\int_0^1\alpha',dt\right)^2
\le
\left(\int_0^1\cos^2\eta,\alpha'^2dt\right)
\left(\int_0^1\sec^2\eta,dt\right).
\]
Therefore the minimum phase energy for a prescribed net change \(a\) is
\[
\frac{a^2}{I_c(\eta)},
\]
attained when
\[
\alpha'(t)=\frac{a,\sec^2\eta(t)}{I_c(\eta)}
\]
with the appropriate sign.

Similarly,
\[
\inf_\beta
\int_0^1\sin^2\eta,\beta'^2dt
=
%
\frac{b^2}{I_s(\eta)}.
\]
If \(I_s=+\infty\), the phase change can be concentrated arbitrarily near a point where \(\eta\) is arbitrarily small, at vanishing energy cost. Geometrically, at \(\eta=0\) the \(w\)-coordinate vanishes and its phase can be reset.

The winding numbers in \(\alpha\) and \(\beta\) are minimized by the principal differences (a,b). This proves \(32\). Allowing the boundary value \(\eta=\pi/4\) does not alter the infimum: a boundary path can be pushed to \(\eta<\pi/4\), fixing its endpoints and changing its length by an arbitrarily small amount. \(\square\)

On any interval on which a minimizing profile satisfies
\[
0<\eta<\frac{\pi}{4},
\]
its Euler--Lagrange equation is
\[
\eta''
+
p^2\frac{\sin\eta}{\cos^3\eta}
------------------------------
%
q^2\frac{\cos\eta}{\sin^3\eta}
=0,
\tag{35}
\]
where
\[
p=\frac{a}{I_c(\eta)},
\qquad
q=\frac{b}{I_s(\eta)}.
\]
At \(\eta=\pi/4\) this is supplemented by the usual obstacle variational inequality; contact with \(\eta=0\) allows a change of the \(\beta\)-phase.

\medskip\hrule\medskip

\subsection{6. Boundary and sanity checks}

\subsubsection{Dimension \(n=1\)}

Here
\[
GL_1^+(\mathbb R)=(0,\infty),
\]
which is convex, so
\[
\boxed{d_{\Omega_1}(a,b)=|a-b|.}
\]
The singular locus is the singleton (\{0\}).

This also confirms that the relevant metric is the induced Euclidean metric, not a left-invariant logarithmic metric.

\subsubsection{Positive scalar multiples}

If \(B=cA\) with \(c>0\), the radial segment is admissible in either cone, and
\[
d(A,B)=|c-1|,|A|_F.
\]

\subsubsection{Equal polar factors}

If
\[
A=QP,\qquad B=QR,
\]
with \(Q\in SO(n)\) and \(P,R>0\), then
\[
(1-t)A+tB=Q((1-t)P+tR)\in GL_n^+.
\]
Hence
\[
d_{\Omega_n}(A,B)=|A-B|_F.
\]

\subsubsection{Equality without an admissible straight segment}

For even \(n\),
\[
A=I,\qquad B=-I
\]
both belong to \(GL_n^+\). The only negative eigenvalue of \(A^{-1}B=-I\) is (-1), with even multiplicity \(n\). Therefore
\[
\boxed{d_{\Omega_n}(I,-I)=2\sqrt n.}
\]
The straight segment passes through \(0\), so no curve in \(GL_n^+\) attains this infimum. Indeed, any Euclidean curve whose length equals its endpoint chord must itself be the straight segment.

This shows that ``geodesic distance'' here is genuinely an infimum and need not be realized.

\subsubsection{Strict inequality in \(GL_2^+\)}

Take
\[
A=I,\qquad B=\operatorname{diag}(-1,-2).
\]
Then
\[
\det((1-t)A+tB)=(1-2t)(1-3t),
\]
which is negative for \(1/3<t<1/2\). Hence the distance is strictly greater than
\[
|A-B|*F=\sqrt{13}.
\]
At \(t=2/5\),
\[
C_t=\operatorname{diag}(1/5,-1/5),
\qquad
\sigma_{\min}(C_t)=\frac15.
\]
Formula \(17\) gives
\[
d_{\Omega_2}(A,B)
\ge
\frac{\sqrt{53}+\sqrt{118}}5
\approx3.62858
%
%
%
\sqrt{13}\approx3.60555.
\]

\subsubsection{A singular pencil without a common kernel}

Let
\[
A=
\begin{pmatrix}
0&0&0\\
0&0&-1\\
0&1&0
\end{pmatrix},
\qquad
B=
\begin{pmatrix}
0&0&1\\
0&0&0\\
-1&0&0
\end{pmatrix}.
\]
Then
\[
\ker A=\mathbb Re_1,\qquad \ker B=\mathbb Re_2,
\]
so there is no common right kernel. Nevertheless every linear combination of \(A\) and \(B\) is a \(3\times3\) skew-symmetric matrix, hence singular. Therefore
\[
d_{\Sigma_3}(A,B)=|A-B|_F.
\]
This shows why the correct condition is singularity of the whole pencil, not existence of a fixed common null vector.

\medskip\hrule\medskip

\subsection{7. Topological and foundational assumptions}

All arguments take place in the finite-dimensional Euclidean space \(M_n(\mathbb R)\).

\begin{itemize}
\item \textbf{Separation and countability.} The spaces are Hausdorff, metrizable, first countable, second countable and separable automatically. No nonmetrizable or net-based argument occurs.
\item \textbf{Path-connectedness.} \(GL_n^+\) is piecewise \(C^1\)-arc-connected by polar decomposition. The singular locus is star-shaped about \(0\), hence path-connected; its punctured unit link is connected for \(n\ge2\).
\item \textbf{Compactness.} No compactness of \(GL_n^+\) is assumed. Compactness is used only in standard finite-dimensional forms, such as compactness of orthogonal groups or unit links. In particular, \(GL_n^+\) is not closed and the infimum need not be attained.
\item \textbf{Sequential versus general topology.} Every limiting argument can be formulated sequentially because all spaces are metric. The main strictness proofs above avoid limiting compactness altogether by using the explicit obstacle bound.
\item \textbf{Regularity of curves.} All constructed curves are piecewise \(C^1\). In the \(2\times2\) variational formula, \(H^1\) or absolutely continuous curves give the same infimum by piecewise-smooth approximation.
\item \textbf{Choice.} No nontrivial choice principle is used; only finite-dimensional singular-value and polar decompositions and finite choices of orthonormal bases occur.
\item \textbf{Norm.} The Hilbertian Frobenius norm is essential to the cosine-law cone formula. A different norm would give a different intrinsic metric.
\item \textbf{Determinant component.} Restricting to \(GL_n^+\) is essential. The two determinant-sign components of \(GL_n(\mathbb R)\) cannot be joined by an invertible path.
\end{itemize}

\medskip\hrule\medskip

\subsection{8. Exact remaining obstruction}

For arbitrary \(n\), Theorem 1 reduces the problem exactly to evaluating
\[
\delta_+ \quad\text{on}\quad
Y_n^+={X\in S_F:\det X>0},
\]
and
\[
\delta_\Sigma \quad\text{on}\quad
Z_n={X\in S_F:\det X=0}.
\]

For \(n=2\), \(Z_2\) is a flat Clifford torus, which gives the closed formula \(28\), while \(Y_2^+\) is a solid Clifford torus and yields the scalar problem \(32\).

For \(n\ge3\), the singular link is stratified by rank:
\[
Z_n=\bigcup_{k=1}^{n-1}{X\in S_F:\operatorname{rank}X=k}.
\]
A shortest path can use lower-rank strata as shortcuts, so solving the ordinary geodesic equations only on the smooth rank-(\(n-1\)) stratum is insufficient. Likewise, paths in \(Y_n^+\) may approach several boundary rank strata.

A relevant result of Katz, Katz, Kerner and Liokumovich proves global bi-Lipschitz comparison between the intrinsic and ambient metrics on \(GL_n^+\), and their rank-stratification argument establishes the corresponding normal-embedding property for the determinantal variety. Its conclusion is quantitative comparability \(d_{\mathrm{int}}\le C(n)d_{\mathrm{ext}}\), rather than an exact distance formula. ([arXiv][2])

The most promising next lemma would be a \textbf{stratified geodesic normal form}: prove that a minimizing curve on the closed unit determinant link can be decomposed into finitely many smooth geodesic pieces lying in fixed-rank strata, with explicitly stated tangent-cone transmission conditions and a number of rank transitions bounded in terms of \(n\). Such a result would convert the remaining angular problem into a finite-dimensional boundary-value optimization.

[1]: https://lucpommeret.com/assets/Qsansreponse260405.pdf "https://lucpommeret.com/assets/Qsansreponse260405.pdf"
[2]: https://arxiv.org/pdf/1602.01227 "https://arxiv.org/pdf/1602.01227"
